\section{Background}

\subsection{Architectural Variants}

There are two well-established architectural families in MAS: Centralized and Decentralized \citep{yang2025agentnetdecentralizedevolutionarycoordination}. Each architectural family can be implemented with varying design choices such as different planning strategies, subagent-autonomy, and context organization that significantly influence the system's vulnerability surface.

\textbf{Centralized Architecture:} A single coordinating agent (often referred to as orchestrator) generally decomposes the user’s request into a plan, assigns the subtasks to specialized agents (also called as subagents), and aggregates the outputs into a final response. In practice, frameworks like OpenAI Agents \citep{OpenAIAgentsSDK2025}, Magentic-One \citep{fourney2024magenticonegeneralistmultiagentsolving}, and LangGraph \citep{LangGraph} provide abstractions for building MAS. They include ready-to-use implementations offering developers convenient starting points for MAS deployment. These implementations vary substantially in their specific design choices even for the same architectural pattern. For instance, two centralized implementations might differ in how the orchestrator delegates tasks, the level of autonomy granted to subagents, or the mechanisms used for aggregating responses, all while maintaining the same fundamental centralized structure. 

\textbf{Decentralized Architecture:} This distributes decision-making authority across agents, typically allowing peer-to-peer communication and local coordination. This reduces reliance on any one of the agents, but raises the risk of mis-coordination, confusion or conflicts. Similar to the case of centralized, the MAS frameworks also offer ready-to-use implementation of this architectural pattern. 

\vspace{-5pt}
\subsection{Adversarial Prompting Benchmarks}
We evaluate MAS on four widely used safety benchmarks. These datasets cover a spectrum of challenges across varying domains like code, web based tasks and file-system based tasks etc. 

\textbf{RedCode~\citep{guo2024redcode}:} is designed to probe whether agents can be guided to generate and execute harmful code when prompted explicitly. This stresses low-level execution vulnerabilities that arise when LLM agents interact directly with system-level commands. 

\textbf{SafeArena~\citep{tur2025safearena}:} tests agent behavior in completing harmful or unethical tasks via interaction with offline replicas of real-world platforms such as Reddit, GitLab, and e-commerce sites. It captures higher-level decision-making vulnerabilities in realistic, environment-driven scenarios. 

\textbf{AgentHarm~\citep{andriushchenko2024agentharm}:} evaluates robustness against jailbreak-style adversarial prompting, where users attempt to coerce agents into unsafe tool use or policy violations. 

\textbf{Agent Security Bench (ASB)~\citep{zhang2024agent}:} evaluates the robustness against Prompt Injection attacks. In this threat model, malicious instructions are added to the benign input prompt, aiming to override the agent's original benign instructions and compel it to perform an unintended or harmful task. This exposes systemic weaknesses in communication and delegation pathways.
   

\subsection{Evaluation Metrics}
Prior works have primarily relied on aggregate outcome measures such as \textbf{Attack Success Rate (ASR)} and \textbf{Refusal Rate (RR)}, which track whether a harmful task was executed or rejected \citep{guo2024redcode, zhang2024agent, andriushchenko2024agentharm}. While useful as high-level summaries, these metrics are limited in two ways: (i) They cannot localize which component of an agentic architecture was responsible for refusal or execution. (ii) They were designed for single-agent settings, making them insufficient for capturing the dynamics of MAS. Recent work by \citet{tur2025safearena} introduced the \textbf{Agent Risk Assessment (ARIA)} framework, that categorizes harmful task outcomes into four levels of risk, reflecting both intent and timing of refusal: 1) ARIA-1 (Immediate Refusal): The agent rejects the task outright. 2) ARIA-2 (Delayed Refusal): The agent begins to comply but ultimately refuses. 3) ARIA-3 (Incorrect Execution): The agent performs harmful actions but fails to complete the task. 4) ARIA-4 (Critical Risk): The agent fully completes the harmful task. ARIA provides a more fine-grained behavioral spectrum than ASR/RR, but it still does not reveal the weak links within the multi-agent pipeline risk, i.e. which specific planner, sub-agent or coordination mechanism introduced or failed to mitigate risk.

\vspace{-10pt}
\section{Challenges}

Finding weak links in MAS faces substantial methodological and technical barriers that prevent systematic security evaluation across different architectures and frameworks.
The key challenges are:


\textbf{Integration Complexity:} MAS evaluation requires integrating existing frameworks like OpenAI Agents, Magentic-One, and LangGraph to work on specialized backends (such as sandboxed environments for RedCode, web browsers for SafeArena, or simulated file-systems) as mandated by the benchmarks. These frameworks substantially differ in code design even when representing the same MAS architecture. Thus, individual development efforts are required to integrate each of the frameworks with every benchmark. As a result, MAS evaluation today is not only slow and inconsistent, but also fundamentally brittle and hinders the community’s ability to obtain reliable and comparable insights into the security and performance of multi-agent architectures.

\textbf{Insufficient Evaluation Metrics:} Existing metrics like Attack Success Rate (ASR), Refusal Rate (RR), and even the more nuanced ARIA framework fail to identify where within MAS, vulnerabilities originate or are mitigated. These outcome-based measures cannot localize whether refusals stem from planning strategies, delegation mechanisms, sub-agent autonomy levels, or communication protocols. This limitation prevents systematic identification of weak links across different design choices within centralized and decentralized architectural families.

\textbf{Lack of Systematic Comparison Framework:} The heterogeneity in framework implementations where different systems vary substantially in design choices while preserving underlying architectural patterns, combined with dataset-specific evaluation requirements, makes reproducible and comparable analysis hard. Without unified evaluation capabilities, researchers cannot systematically compare different MAS implementations to identify common vulnerability patterns or validate whether security insights generalize across architectural variants and application domains.